\documentclass[9pt,twocolumn,twoside]{article}

\usepackage[T1]{fontenc}
\usepackage{newtxtext,newtxmath}
\usepackage{helvet}
\usepackage[margin=1.9cm,top=2.2cm,bottom=2.4cm]{geometry}
\usepackage{titlesec}
\usepackage{booktabs}
\usepackage{fancyhdr}
\usepackage{caption}
\usepackage{enumitem}
\usepackage[colorlinks=true,allcolors=blue]{hyperref}
\usepackage{microtype}

\setlength{\columnsep}{0.8cm}

% --- OUP-style section headings: bold sans, unnumbered -----------------
\titleformat{\section}{\normalfont\sffamily\bfseries\large}{}{0pt}{}
\titlespacing{\section}{0pt}{1.1em}{0.4em}
\titleformat{\subsection}{\normalfont\sffamily\bfseries\normalsize}{}{0pt}{}
\titlespacing{\subsection}{0pt}{0.9em}{0.3em}
\setcounter{secnumdepth}{0}

% --- captions: "Table 1." bold, above tables ---------------------------
\captionsetup[table]{labelfont=bf,labelsep=period,justification=raggedright,
  singlelinecheck=false,position=above,font=small,skip=4pt}

% --- headers/footers ----------------------------------------------------
\pagestyle{fancy}
\fancyhf{}
\fancyhead[LE]{\small\thepage\quad\itshape B.\,Judy}
\fancyhead[RO]{\small\itshape Storm-Level Verification of Operational Cloud Seeding\quad\thepage}
\fancyfoot[C]{}
\renewcommand{\headrulewidth}{0pt}

\fancypagestyle{firstpage}{%
  \fancyhf{}
  \fancyhead[L]{\footnotesize Working paper, v1.1\\First version: 20 July 2026; this version: 21 September 2026\\Code, data and figures: \url{https://github.com/bryjudy/seeding-verification}\\}
  \fancyfoot[L]{\footnotesize \copyright\ Bryce Judy 2026. Data from public U.S.\ federal and Utah state sources. For permissions, please e-mail: bryce.judy@tarifflo.com}
  \fancyfoot[R]{\footnotesize\thepage}
  \renewcommand{\headrulewidth}{0.8pt}
  \renewcommand{\footrulewidth}{0.4pt}
}

\newcommand{\keywords}[1]{\vspace{0.5em}\noindent\textsf{\textbf{Key words:}} #1}

% numeric citations, hand-numbered to match the reference list
\newcommand{\citeparen}[1]{[\csname refnum#1\endcsname]}
\expandafter\def\csname refnumfrench2018\endcsname{1}
\expandafter\def\csname refnumfriedrich2020\endcsname{2}
\expandafter\def\csname refnumgriffith2009\endcsname{3}
\expandafter\def\csname refnumdwr\endcsname{4}
\expandafter\def\csname refnumhrrr\endcsname{5}
\expandafter\def\csname refnumawdb\endcsname{6}
\expandafter\def\csname refnumborusyak2024\endcsname{7}

\begin{document}

\twocolumn[{%
\thispagestyle{firstpage}
\vspace*{0.8cm}
{\noindent\sffamily\bfseries\LARGE Storm-Level Independent Verification of\\[2pt] Operational Cloud Seeding: Evidence from\\[2pt] Utah's 2024--25 Winter Season\par}
\vspace{0.9em}
{\noindent\sffamily\bfseries\large Bryce Judy\textsuperscript{1,$*$}\par}
\vspace{0.6em}
{\noindent\footnotesize $^{*}$Corresponding author. \href{mailto:bryce.judy@tarifflo.com}{bryce.judy@tarifflo.com}\par}
\vspace{1.2em}
{\noindent\sffamily\bfseries\large Abstract\par}
\vspace{0.35em}
{\noindent\small
Operational winter cloud seeding programs claim precipitation increases of
3--15\%, but the evaluations supporting these claims are produced by program
operators using target/control methods that cannot separate seeding effects
from \emph{operator storm selection}: storms are seeded precisely when
atmospheric conditions favor the target areas. This study performs, to the
author's knowledge, the first independent storm-level evaluation of an
operational program built entirely from public data. We assemble a
machine-readable record of Utah's 2024--25 operations---181 generator sites,
203 seeded storm periods totaling 17{,}143 generator-hours across seven
programs---together with per-storm 700\,hPa winds from the HRRR archive and
daily precipitation from 546 SNOTEL stations in seven states. Three research
designs of increasing selection-robustness are applied. A naive seeded-versus-unseeded
storm comparison yields an implausible $+335\%$ for the largest program,
quantifying the selection artifact. Same-sector near-miss controls reduce the
pooled estimate to $+24\%$ [$+4$, $+50$]. The most robust design---a
wind-rotation test in which plume exposure is recomputed from each storm's own
wind field per generator---yields $+1.5\%$ [$-9.6$, $+13.1$] per seeded storm
with a clean placebo on unseeded storms. One season therefore cannot
distinguish the claimed 3--15\% effect from zero, while per-storm effects
larger than $\sim$15--20\% are excluded. The design's power scales directly
with additional seasons of operations records. All data and code are released.
\par}
\vspace{0.6em}
{\noindent\small\keywords{Weather Modification, Cloud Seeding, Causal Inference, SNOTEL, HRRR, Program Evaluation}\par}
\vspace{1.6em}
}]

\section{Introduction}

Winter orographic cloud seeding disperses silver iodide (AgI) into
supercooled liquid water within storms passing over mountain barriers,
nucleating ice crystals that fall as additional snow. The physical mechanism
is established: the SNOWIE field campaign traced seeding material through
individual storms with airborne radar and observed unambiguous seeding
signatures \citeparen{french2018}, and subsequent analysis quantified the
snowfall attributable to individual seeding events \citeparen{friedrich2020}.
What remains contested is the \emph{operational} question: how much water do
routine, ground-based seeding programs actually add, season over season?

The stakes of that question are growing. Utah operates the largest
remotely-controlled seeding generator network in the United States,
expanded in 2023 with \$12M in one-time funding and a \$5M annual budget
following record-low levels in the Great Salt Lake \citeparen{dwr}. The
State's published evaluations report target-area increases of 5--20\% and an
estimated 249{,}600 acre-feet of average annual augmented runoff at
approximately \$1 per acre-foot \citeparen{griffith2009}. Colorado River
basin states co-fund seeding in Utah under multi-state agreements, meaning
these estimates influence interstate water accounting.

These evaluations, however, share two structural weaknesses. First, they are
produced by the contractor operating the programs. Second, and more
fundamentally, they rest on historical target/control regressions whose
treatment variable is not exogenous: \emph{operators seed exactly those
storms whose wind direction, moisture, and temperature favor the target
areas} \citeparen{griffith2009}. Suspension rules compound the problem ---
Utah's criteria suspend seeding when snowpack exceeds 150--200\% of average
\citeparen{griffith2009}, making treatment anti-correlated with natural
snowfall at seasonal scale and positively selected at storm scale. Any
comparison of seeded versus unseeded periods, at any aggregation, inherits
this selection unless the design explicitly breaks it.

This paper contributes three things. (1)~A machine-readable reconstruction
of one full season of Utah seeding operations from public documents:
generator coordinates, per-storm seeding dates and generator-hours for all
seven ground-based programs. (2)~A sequence of three storm-level research
designs applied to the same season of SNOTEL precipitation data, constructed
so that each successive design removes a channel of selection bias; the
progression of estimates ($+335\% \rightarrow +24\% \rightarrow +1.5\%$)
directly measures how much of a naive ``seeding effect'' is selection
artifact. (3)~A wind-rotation identification strategy, which to the author's
knowledge has not previously been applied to an operational program: because
AgI plumes travel downwind, a station's exposure changes storm-by-storm with
the storm's own wind direction, and a real seeding effect must rotate with
it --- something neither operator storm selection nor static orography can
mimic.

\section{Data}

All inputs are public; acquisition and processing code is released with the
paper.

\subsection{Seeding operations records}
The Utah Division of Water Resources publishes annual reports for each
seeding program, prepared by the operator (North American Weather
Consultants) \citeparen{dwr}. From the eight 2024--25 reports we extract:
(i)~generator site tables --- 181 sites with coordinates across the
Central/Southern Utah + Tooele, Western Uintas, High Uintas, Northern Utah,
East Shore, Six Creeks, and Cache Valley programs (Book Cliffs sites are
published as town names only and are geocoded approximately); and
(ii)~storm-operations tables --- 203 program-storm periods with dates, sites
used, and generator-hours, totaling 17{,}143 hours across seven ground-based
programs. Program treatment histories (adoption years, the 1983--87
statewide suspension) are taken from the operator's 30-season retrospective
\citeparen{griffith2009}.

\subsection{Storm winds}
For every storm period and detected precipitation event we compute the mean
700\,hPa wind at each program's generator network from the NOAA
High-Resolution Rapid Refresh (HRRR) analysis archive on AWS Open Data
\citeparen{hrrr}, sampling 6-hourly analyses and extracting only the two
relevant GRIB2 messages per file via byte-range requests. Storm-mean winds
are vector-averaged. Consistent with program design assumptions, 95\% of
seeded storms have 700\,hPa flow from the WSW--NW quadrant (median
264$^\circ$).

\subsection{Snowpack and precipitation observations}
Daily accumulated-precipitation (PREC) series for all 546 SNOTEL stations in
Utah, Colorado, Nevada, Wyoming, Idaho, Arizona, and New Mexico are drawn
from the USDA NRCS Air-Water Database \citeparen{awdb}. Daily increments are
computed by first-differencing, with negative and physically implausible
($>$8\,in/day) increments removed. A 45-year station-winter panel
(1981--2026; 18{,}067 station-winters) is used for the seasonal-scale
benchmark design.

\section{Methods}

\subsection{Design 0: seasonal panel benchmark}
As a baseline we estimate a staggered-adoption imputation estimator
\citeparen{borusyak2024} on log November--March precipitation over 45
winters: station and winter fixed effects are fit on untreated observations
(out-of-state controls, excluding ranges with their own seeding programs,
plus pre-adoption years of Utah cohorts), and treatment effects are read off
as anomalies of treated station-winters. Suspension years are excluded from
fitting because they are endogenously wet.

\subsection{Design 1: seeded vs.\ unseeded storms, far controls}
Storm events are detected as contiguous days of regional precipitation.
For each program, each event is labeled seeded/unseeded from the operations
record; the outcome is the log ratio of mean storm-total precipitation at
plume-exposed stations ($\leq$40\,km downwind of a generator) to that at far
control stations ($>$90\,km from all generators). The effect is the
difference in mean log ratio between seeded and unseeded storms. This is
essentially the operational-evaluation logic transposed to storm scale, and
it is deliberately included as a measurement of selection bias.

\subsection{Design 2: near-miss controls}
Identical, except controls are ``near-miss'' stations 55--130\,km from the
program's generators \emph{in the same downwind sector}: they experience the
same storm types that trigger seeding decisions but lie beyond AgI transport
range. Storm-type selection therefore cancels to first order in the
seeded-minus-unseeded contrast. Inference is by bootstrap over storm events.
Two falsification tests are run: a placebo in which all seeded windows are
shifted by $+10$ days, and a dose-response test correlating per-storm
generator-hours with the per-storm effect.

\subsection{Design 3: wind-rotation (per-generator plume geometry)}
The most selection-robust design exploits within-station, across-storm
variation in exposure. For each station within 70\,km of a program's
generators and each detected event $e$, the station is \emph{in-plume} if any
generator lies 5--45\,km upwind such that the bearing from generator to
station is within $\pm 40^\circ$ of event $e$'s downwind direction (from the
HRRR storm-mean wind). We estimate
\[
\log y_{se} = \alpha_s + \gamma_e + \delta\,\mathrm{score}_{se} +
\beta_1 P_{se} S_e + \beta_0 P_{se}(1-S_e) + \varepsilon_{se},
\]
where $\alpha_s, \gamma_e$ are station and program-event fixed effects,
$\mathrm{score}_{se} = \max_g \cos(\Delta\theta_{sg,e})\,e^{-d_{sg}/40}$ is a
smooth wind-alignment control absorbing orographic anisotropy that rotates
with the wind, $P_{se}$ is the in-plume indicator, and $S_e$ indicates seeded
events. $\beta_1$ is the seeding effect; $\beta_0$ is a built-in placebo that
must be zero, since no AgI is released in unseeded events. Fixed effects are
removed by iterated two-way demeaning; confidence intervals are by bootstrap
over events. Selection cannot generate $\beta_1 > 0 = \beta_0$: operators
choose \emph{when} to seed, but they cannot rotate a precipitation
enhancement around their generators in step with each storm's wind.

\section{Results and Analysis}

\subsection{Seasonal benchmark: the claimed effect is below the noise floor}
The 45-year panel estimator yields an overall treatment effect of $-1.4\%$
(95\% CI $[-6.0, +3.6]$) on winter precipitation across all Utah adoption
cohorts, with a pre-adoption placebo of $+5.6\%$. Seasonal target/control
analysis of station data therefore has a noise floor of roughly $\pm 6\%$
--- wider than the lower half of the claimed 3--15\% range. Published
seasonal evaluations reporting significant 5--20\% increases from comparable
data should be interpreted with this in mind.

\subsection{Selection bias is enormous at storm scale}
Design 1 produces a pooled ``effect'' of $+69\%$, and $+335\%$ for the
Central/Southern program (Table~\ref{tab:designs}) --- an order of magnitude
beyond physical plausibility \citeparen{friedrich2020} and a direct
quantification of operator storm selection: seeded storms are storms chosen
for favoring the targets.

\begin{table}[t]
\caption{Pooled storm-level estimates across designs, winter 2024--25.
Effects are percentage changes in storm-total precipitation at exposed
stations; 95\% CIs by bootstrap over storm events.}
\label{tab:designs}
\small
\begin{tabular}{@{}lrr@{}}
\toprule
Design & Effect & 95\% CI \\
\midrule
1. Far controls (naive) & $+69.0\%$ & $[-9.9, +220.8]$ \\
2. Near-miss controls & $+24.2\%$ & $[+3.7, +49.7]$ \\
\quad placebo: windows $+10$ days & $+0.2\%$ & $[-16.7, +20.7]$ \\
3. Wind-rotation, per-generator & $+1.5\%$ & $[-9.6, +13.1]$ \\
\quad placebo: unseeded events & $-8.4\%$ & $[-26.8, +10.3]$ \\
\bottomrule
\end{tabular}

\medskip
\footnotesize Source: author's calculations from DWR operations reports,
HRRR, and SNOTEL data; full per-program results in the released repository.
\end{table}

\subsection{Near-miss controls: directionally positive, still inflated}
Design 2 collapses the pooled estimate to $+24.2\%$ [$+3.7$, $+49.7$] over
121 program-storm observations. The $+10$-day placebo returns $+0.2\%$
[$-16.7$, $+20.7$], confirming the pipeline produces no mechanical effect;
per-storm generator-hours correlate positively but insignificantly with the
per-storm effect (Spearman $\rho = 0.13$, $p = 0.26$, $n=75$). The estimate
nevertheless remains above the physically plausible range, consistent with
residual selection: orographic enhancement on strong-flow seedable storms is
naturally larger at the barrier than 55--130\,km downwind, and near-miss
controls only partially absorb this.

\subsection{Wind-rotation: no resolvable effect at one season}
Design 3 yields $\beta_1 = +1.5\%$ [$-9.6$, $+13.1$] for in-plume exposure
during seeded events, with the unseeded-event placebo at $\beta_0 = -8.4\%$
[$-26.8$, $+10.3$] --- statistically zero, as required. (An intermediate
specification using program-centroid rather than per-generator geometry
produced an unseeded-event artifact of $+12\%$, illustrating that plume
geometry matters: several Utah generator networks form north--south lines
exceeding 100\,km, which a centroid represents poorly.)

The point estimate sits exactly where the physical literature suggests a
true operational effect should --- small and positive --- but one season of
data cannot distinguish it from zero. Conversely, per-storm effects larger
than approximately 15--20\% are excluded at conventional confidence levels.

\subsection{Power and the path to resolution}
The binding constraint is seasons, not stations. Winter 2024--25 supplies 26
usable events and 121 program-event observations, giving Design 3 a CI
half-width of $\sim$11 percentage points. Under simple $1/\sqrt{n}$ scaling,
four to nine additional seasons of operations records --- which the Division
of Water Resources maintains and provides on request --- would narrow the CI
to $\pm 4$--$5\%$, sufficient to confirm or refute the 3--15\% operational
claims. The marginal cost of extending the analysis is transcription of the
operators' storm tables; all other inputs (HRRR, SNOTEL) extend
automatically.

\section{Discussion}

Three conclusions follow. First, the headline numbers used to justify
operational seeding budgets are produced by methods whose selection bias, at
storm scale, we measure at more than an order of magnitude larger than any
plausible true effect; seasonal versions of those methods have noise floors
wider than the effects claimed. Second, the most selection-robust design
available finds a small positive point estimate fully consistent with both
the gold-standard physical evidence \citeparen{friedrich2020} and with zero
--- after one season. The correct current statement about Utah's program is
therefore agnostic: the claimed effect is neither confirmed nor refuted, and
effects large enough to matter for interstate water accounting
($\gtrsim$15--20\% per storm) are ruled out. Third, extra-area
(``downwind'') effects on neighboring regions --- a recurring political
concern --- are necessarily smaller than target-area effects and hence well
below ground-based detectability at current program intensities.

\subsection{Limitations}
Wind sectors use each storm's mean 700\,hPa flow at the network; hourly
wind shifts within storms are not modeled. SNOTEL gauges undercatch snowfall
in wind, which attenuates but does not bias the estimates given fixed
effects. Book Cliffs generator locations are approximate. The seeded/unseeded
event labels inherit any errors in the operators' storm tables. Finally, the
event catalog is built from precipitation, so completely dry seeded periods
are unobservable; this biases against detecting effects only if seeding
generates precipitation where none would otherwise occur, which the physical
literature does not support at meaningful scale.

\section{Data and Code Availability}
All datasets constructed for this study (generator coordinates, storm
operations record, storm winds, event-level results) and all analysis code
are released at \url{https://github.com/bryjudy/seeding-verification}. Every input derives from public
sources: Utah DWR seasonal reports, NOAA HRRR (AWS Open Data), and USDA NRCS
SNOTEL.

\section{Acknowledgments}
The author thanks the Utah Division of Water Resources for publishing
operational seasonal reports at the level of detail that makes independent
evaluation possible, and acknowledges the use of Claude (Anthropic) for
research and engineering assistance throughout.

\section{References}
\newcounter{refc}
\newcommand{\bibitemx}{\stepcounter{refc}\item[\arabic{refc}.]}
\begin{itemize}[leftmargin=1.6em,itemsep=2pt,parsep=0pt]\small
\bibitemx French, J.~R., et al. ``Precipitation formation from orographic
cloud seeding.'' \emph{Proceedings of the National Academy of Sciences}
115.6 (2018): 1168--1173.
\bibitemx Friedrich, K., et al. ``Quantifying snowfall from orographic cloud
seeding.'' \emph{Proceedings of the National Academy of Sciences} 117.10
(2020): 5190--5195.
\bibitemx Griffith, D.~A., M.~E. Solak, and D.~P. Yorty. ``30+ Winter Seasons
of Operational Cloud Seeding in Utah.'' \emph{Journal of Weather
Modification} 41 (2009). AMS preprint:
\url{https://ams.confex.com/ams/pdfpapers/138877.pdf}.
\bibitemx Utah Division of Water Resources. ``Cloud Seeding'' program page
and 2024--25 seasonal reports. Accessed July 2026.
\url{https://water.utah.gov/cloudseeding/}.
\bibitemx NOAA. ``High-Resolution Rapid Refresh (HRRR).'' Registry of Open
Data on AWS. \url{https://registry.opendata.aws/noaa-hrrr-pds/}.
\bibitemx USDA Natural Resources Conservation Service. ``Air and Water
Database (AWDB) REST API / SNOTEL.''
\url{https://wcc.sc.egov.usda.gov/awdbRestApi/}.
\bibitemx Borusyak, K., X.~Jaravel, and J.~Spiess. ``Revisiting Event-Study
Designs: Robust and Efficient Estimation.'' \emph{Review of Economic
Studies} 91.6 (2024): 3253--3285.
\end{itemize}

\end{document}
